% Round 20, family20 branch. Candidate for exact independent review.
\section{A surface incidence inverse and arbitrary good-reduction coefficients}
\label{r20as:section}

\paragraph{Scope.}
The result of this section concerns \emph{all} mixed homological cycles
on $T\times S$, where $T$ has potentially good projective reduction and
$S$ is an arbitrary smooth projective surface.  Neither a Chow
decomposition of $T$ nor a numerical--homological comparison on its
reduction is required.  The all-degree abelian, supplied abelian-motive,
and good-factor-times-abelian cases were already frozen in R16 and
are not counted as new here.  The same good-coefficient-times-surface
scope was proposed in R9; the present section gives an independent
proof and verifies it through the current AK19/CR19/VT17 interfaces.
It is not presented as a new family idea.  The source of the new proof is a
correspondence identity on the \emph{special-fiber surfaces}; it is
not a proposed arbitrary proper pushforward of actual forms.

All Chow groups have rational coefficients.  Numerical equivalence
and degrees on special fibers are native over their finite residue
field.  Cohomology means geometric $\ell$-adic cohomology, with the
indicated Tate twists, and $\ell$ is different from the residue
characteristic.

\subsection{The actual incidence correspondence}

Let $R$ be a complete discrete valuation ring with finite residue
field $k$, and let $\mathcal S/R$ be a regular projective reduced
strictly semistable model of a smooth projective surface.  After a
finite unramified extension, we may and do assume that the finitely
many special-fiber components $D_i$ are geometrically integral.
Write $D=\coprod_iD_i$.  Each $D_i$ is a smooth projective surface.
For each double component $E_e\subset D_i\cap D_j$, let
$j_{e,i}:E_e\hookrightarrow D_i$ and $j_{e,j}:E_e\hookrightarrow D_j$
be its two embeddings.  Orienting $e$ from $i$ to $j$, use the
frozen signs
\[
 \rho(x)_e=j_{e,j}^*x_j-j_{e,i}^*x_i,
 \qquad
 \sigma(y)_i=\sum_{e\text{ out of }i}j_{e,i*}y_e
              -\sum_{e\text{ into }i}j_{e,i*}y_e.
 \tag{AS20.1}
\]
Thus $A=\sigma\rho$ is represented by an actual rational
correspondence
\[
             \mathcal A\in\operatorname{CH}^{3}(D\times_kD).
 \tag{AS20.2}\label{r20as:Acorrespondence}
\]
Its entry from $D_i$ to a different $D_j$ is
$\sum_{e:i\leftrightarrow j}(j_{e,i},j_{e,j})_*[E_e]$.
Its diagonal entry is
\[
       -\Delta_{D_i*}\left(\sum_{e\text{ incident to }i}[E_e]\right).
 \tag{AS20.3}
\]
The principal total special fiber gives the equivalent normal-bundle
description of the diagonal.  The formulas hold on actual Chow
groups, on their numerical quotients, and on full geometric
cohomology.  They do not delete triple points.  A double curve may
meet other double curves; these are precisely the permitted higher
strata of a surface model.

\begin{lemma}[Numerical generalized inverse as a correspondence]
\label{r20as:numerical-inverse}
There is an actual rational divisor correspondence
\[
             \mathcal G\in\operatorname{CH}^{1}(D\times_kD)
\]
for which
\[
          \mathcal A\circ\mathcal G\circ\mathcal A-\mathcal A
                    \quad\text{is numerically zero}
 \tag{AS20.4}\label{r20as:numerical-identity}
\]
as a correspondence, entry by entry.  The assertion is about
correspondence classes, not merely their action on numerical
cycles of the individual components.
\end{lemma}

\begin{proof}
Use Jannsen's semisimplicity theorem for rational numerical motives
over $k$ \cite[Theorem~1]{AS20Jannsen}.  With the convention that a
degree-one correspondence is a morphism $M\to M(1)$, the numerical
class of $\mathcal A$ is such a morphism for
$M=\bigoplus_i h(D_i)$.  In a semisimple abelian category, split its
source as its kernel plus a complement, and its target as its image
plus a complement.  Invert the resulting isomorphism from the source
complement to the image, and extend this inverse by zero on the target
complement.  The resulting numerical morphism $G_N$ satisfies
$A_NG_NA_N=A_N$.

A morphism of degree minus one between these surface motives is a
numerical class of a divisor on $D_i\times D_j$.  Each such numerical
class has an actual rational Chow representative by definition of
the numerical quotient.  Choosing representatives of the finitely
many entries of $G_N$ gives $\mathcal G$.  Their composition has
exactly the claimed numerical class.  No cycle-class map from a
numerical group is used.
\end{proof}

\begin{proposition}[One rational inverse on full component cohomology]
\label{r20as:cohomological-inverse}
The correspondence $\mathcal G$ of
Lemma~\ref{r20as:numerical-inverse} satisfies
\[
                   A^H G^H A^H=A^H
 \tag{AS20.5}\label{r20as:full-inverse}
\]
on all geometric cohomology of $D$, simultaneously for every
$\ell\ne\operatorname{char}k$.  Its numerical action satisfies the
same identity in every degree.  In particular this is an inverse
on the entire cohomological image of $A$, including odd-degree
cohomology; no algebraicity assertion about all of $H^2(D_i)$ is
made.
\end{proposition}

\begin{proof}
Each entry of the error in \eqref{r20as:numerical-identity} belongs
to $\operatorname{CH}^3(D_i\times D_j)$ and is numerically zero.
The product $D_i\times D_j$ has dimension four, so this error is
an actual \emph{one-cycle}.  By the curve-comparison lemma of SP19,
a numerically zero native one-cycle on a product of two smooth
projective surfaces has zero geometric $\ell$-adic class for every
invertible $\ell$.

For completeness, that comparison uses the known algebraicity of
the inverse of $L^2:H^2\to H^6$ on a product of surfaces.  The
Lefschetz standard conjecture is known for surfaces and is stable
under products \cite[Proposition~4.3]{AS20Kleiman}.  Applying the
inverse correspondence to a numerical-zero curve gives a divisor;
the polarized divisor-index pairing shows that this divisor is
numerically zero, hence geometrically homologically zero.  Applying
$L^2$ again kills the original curve class.  Native numerical zero
continues to pair to zero against geometric divisors after finite
constant extension, by the full-degree norm and projection formula.

Consequently the error correspondence has zero class, rather than
only zero action on an algebraic subspace.  Compatibility of
correspondence composition with geometric cycle classes gives
\eqref{r20as:full-inverse} on every cohomological degree.
The numerical identity is already \eqref{r20as:numerical-identity}.
\end{proof}

\paragraph{A finite rational certificate.}
After choosing rational divisor representatives of bases of each
$N^1(D_i\times D_j)$, write the entries of $\mathcal G$ as unknown
rational combinations of these representatives.  Test every entry
of $\mathcal A\mathcal G\mathcal A-\mathcal A$ against a basis of
the corresponding numerical divisor space.  The equations
\[
 \deg_k\left(
   (\mathcal A\mathcal G\mathcal A-\mathcal A)_{ji}
                  \cdot H_{ji,\mu}\right)=0
 \tag{AS20.6}
\]
form a finite \emph{linear} rational system in the coefficients of
$\mathcal G$: $\mathcal A$ is fixed.  Its solvability follows from
Lemma~\ref{r20as:numerical-inverse}, and the perfect numerical
curve--divisor pairing makes its vanishing equivalent to
\eqref{r20as:numerical-identity}.  All degrees in this certificate
are on proper finite-field fourfolds.  This is a finite certificate
once complete bases are supplied, not a claimed algorithm for
finding a complete Neron--Severi basis on every input.

\subsection{Arbitrary good-reduction coefficients, including mixed classes}

\begin{proposition}[Tensoring the inverse with a good coefficient]
\label{r20as:good-coefficient}
Let $\mathcal T/R$ be any smooth projective model with geometrically
integral generic fiber.  Then
$\mathcal X=\mathcal T\times_R\mathcal S$ is regular projective
strictly semistable.  On its components $T_0\times D_i$, the
AK19 lifting defect and the CR19 supported defect vanish in every
degree:
\[
                 \Delta_{A,a,\ell}(\mathcal X)=0,
       \qquad \mathscr R_{A,r,\ell}(\mathcal X)=0.
 \tag{AS20.7}\label{r20as:coefficient-defects}
\]
There is no condition on numerical versus homological equivalence
on $T_0$ or on $T_0\times D_i$.
\end{proposition}

\begin{proof}
The projection of $\mathcal X$ to $\mathcal S$ is smooth, so
regularity and strict semistability follow by smooth base change.
Its double embeddings and normal line bundles are the products of
the original ones with $T_0$.  Under the reordering of the four
factors, its actual incidence correspondence and the chosen inverse
are
\[
        \mathcal A'=\Delta_{T_0}\boxtimes\mathcal A,
        \qquad
        \mathcal G'=\Delta_{T_0}\boxtimes\mathcal G.
 \tag{AS20.8}
\]
External products commute with composition of correspondences.
The cycle class of
$\mathcal A'\mathcal G'\mathcal A'-\mathcal A'$
is zero by Proposition~\ref{r20as:cohomological-inverse}; equivalently,
\[
                    A'{}^H G'{}^H A'{}^H=A'{}^H
 \tag{AS20.9}\label{r20as:product-inverse}
\]
on \emph{all} cohomology of the product components.  This includes
the mixed Kunneth summands.  No claim that actual Chow classes on
the product are external products is required.

Let $p_u:\mathscr A_C^u\to C^u\otimes\mathbf Q_\ell$ be AK19's
canonical projection from the actual algebraic class span to native
numerical classes.  A genuine Chow correspondence preserves
numerical equivalence, by the correspondence projection formula.
It therefore commutes with $p_u$ in its stated degree, including
for $\mathcal G'$.

For $v\in\ker A'{}^N_a$, choose any actual algebraic cohomology
lift $x\in\mathscr A_C^a$ with $p_a(x)=v$.  The correction
\[
                         x^{\mathrm{corr}}=x-G'{}^H A'{}^H x
 \tag{AS20.10}\label{r20as:actual-correction}
\]
is again in that actual algebraic span.  Its numerical image is
$v-G'{}^N A'{}^N v=v$, and its cohomological $A'$-image is zero
by \eqref{r20as:product-inverse}.  Thus every numerical $A'$-kernel
test has an actual algebraic $A'$-kernel lift.  This proves
$\Delta_A=0$ by AK19's exact connecting description.

One may also check the smaller defect directly.  If
$\alpha\in\mathscr A_C^r\cap\operatorname{im}A'{}^H_{r-1}$,
then $G'{}^H\alpha$ is actual algebraic and
$A'{}^H G'{}^H\alpha=\alpha$.  Applying $p_r$ puts its numerical
class in $\operatorname{im}A'{}^N_{r-1}$, proving
$\mathscr R_A=0$.  This argument uses actual correspondences on
the finite-field components only; no actual analytic form is pushed
along an arbitrary proper map.
\end{proof}

\begin{theorem}[A potentially good factor times any surface]
\label{r20as:full-G}
Fix a finite valued extension $K/\mathbf Q_p$ and its absolute
value.  Let $T/K$ be smooth projective geometrically integral and
suppose that after some finite extension it has a \emph{smooth
projective} model over the valuation ring.  Let $S/K$ be an arbitrary
smooth projective geometrically integral surface.  Put
$n=\dim(T\times S)$.  Then, for every $1\le r\le n$, every
geometrically homologically trivial rational codimension-$r$
Chow class on $T\times S$ has zero actual obstruction.
Consequently every specified such cycle has an ordinary-actual
strong Green form on the complement of its exact original support,
with the identity against every compactly supported actual test.

The conclusion holds afresh for every finite valued extension of
$K$ and every new cycle over that extension.
\end{theorem}

\begin{proof}
If $\dim T=0$, geometric integrality makes $T=\operatorname{Spec}K$,
and the assertion is the frozen surface endpoint theorem.  Assume
henceforth $\dim T\ge1$, so the product dimension is at least three,
as required by the AK19/CR19 geometric application.

Choose first a finite extension on which $T$ has its smooth
projective model.  Apply the frozen projective semistable alteration
theorem to $S$ over that field.  After a further finite extension
$L/K$, it supplies a smooth projective geometrically integral surface
$S'/L$, a dominant generically finite morphism $S'\to S_L$, and a
regular projective strictly semistable model $\mathcal S'/L^\circ$.
The smooth projective model of $T$ remains smooth projective under
every further finite base change.  A further unramified extension
splits the finitely many component constants of $\mathcal S'$.
The product model for $Y=T_L\times S'$ therefore has the vanishing
defects of Proposition~\ref{r20as:good-coefficient}.

All degrees of geometric cohomology of a surface satisfy
weight--monodromy: the known degrees at most two and Poincare
duality cover all degrees; see the recorded known cases following
Conjecture~0.5 in \cite{AS20Saito}.  For $T_L$ the smooth proper
model gives pure cohomology with trivial monodromy.  Kunneth and
the tensor rule for the monodromy filtration thus give
weight--monodromy for every degree of $Y$, in particular for each
$H^{2r-1}(Y_{\bar L},\mathbf Q_\ell(r))$ needed here.  Recovered
OW19 makes the actual inertia term zero.  Together with
\eqref{r20as:coefficient-defects}, AK19 (or CR19) makes every
homological cycle obstruction on $Y$ zero.

For a newly specified class $z$ on $T\times S$, its pullback along
$g:Y\to(T\times S)_L$ is homological, so has zero obstruction.
Frozen R16 dominant Bott--Chern detection over $L$ first gives zero
obstruction for $z_L$.  Finite-constant descent along
$(T\times S)_L\to T\times S$, with the full degree $[L:K]$ and
the extending original absolute value, gives zero obstruction over
the original valued $K$.  VT17 then reconstructs the Green form for
the \emph{originally specified cycle}.  The alteration is not
assumed finite on the support complement.  There is no arbitrary
proper actual-form pushforward and no restriction to exterior-product
tests.

For an arbitrary finite valued extension $K'/K$, the variety $T_{K'}$
still has potentially good projective reduction: base change a
chosen model to a compositum.  The variety $S_{K'}$ is again an
arbitrary smooth projective geometrically integral surface.  Apply
the same argument, including a fresh surface alteration if needed,
to every newly given class.  No assertion that the class descends
from $K$ is used.
\end{proof}

\begin{corollary}[Full complementary real heights]
\label{r20as:height}
Let $F$ be a number field.  Let $T/F$ be smooth projective
geometrically integral with potentially good projective reduction
at every finite place, and let $S/F$ be any smooth projective
geometrically integral surface.  Put $X=T\times S$ and
$n=\dim X$.  The frozen PH16 canonical real pairing has the full
domains
\[
 \operatorname{CH}^r(X)_{\mathbf Q,\mathrm{geom\,hom}}
 \times
 \operatorname{CH}^{n+1-r}(X)_{\mathbf Q,\mathrm{geom\,hom}}
                    \longrightarrow\mathbf R
 \tag{AS20.11}\label{r20as:full-height}
\]
for every complementary positive degree.  The same conclusion
holds over every finite number-field extension for all new cycles.
In particular, when $\dim T=3$, this includes the complete middle
self-pairing on $\operatorname{CH}^3(T\times S)_{\mathbf Q,\mathrm{geom\,hom}}$.
\end{corollary}

\begin{proof}
At every finite completion apply Theorem~\ref{r20as:full-G} in
both complementary degrees.  The complete local homological groups
lie in the exact FD17/VT17 obstruction kernels.  Frozen PH16 and
its archimedean construction give \eqref{r20as:full-height}; GA17
supplies admissible Chow lifts on the specified regular arithmetic
models or alteration witnesses.  After a finite number-field
extension, potential good projective reduction persists at every
finite completion, and the theorem is reapplied to all new classes.
No positivity, nondegeneracy, or $L$-value identity is asserted.
\end{proof}

\paragraph{Further closure with an abelian factor.}
For an abelian variety $B/F$ and an arbitrary surface $S/F$, the
same full height conclusion follows if at each finite place at
least one of the following holds: $B$ has potentially good
reduction, or $S$ has potentially good \emph{projective} reduction.
The first alternative is the new theorem with $T=B$; the second
is the frozen R16 good-factor-times-abelian theorem.  This is a
place-by-place assertion, so no global common local model is being
postulated.  The corresponding finite-extension quantifier follows
by reapplication at the new completions.

\paragraph{Limits and the next product gate.}
The argument proves no general vanishing theorem for an arbitrary
variety.  For a $d$-dimensional bad factor, the incidence error would
be a cycle of dimension $d-1$ on $D_i\times D_j$; only for $d=2$
is this the curve-comparison case just used.  Further, a product of
two bad semistable models is singular before resolution.  Its
resolved component incidence operator is not asserted to be the
naive tensor sum of the two original $A$ operators.  A compatible
splitting of the full stratum system would be additional input.
The present theorem therefore does not, by iteration alone, give
all-degree G on arbitrary products of several bad surfaces.

\paragraph{Sources and historical scope.}
The primary numerical-motive input was read in
\cite[Theorem~1]{AS20Jannsen}.  Curve comparison and its precise
Lefschetz input are the frozen SP19 proof and its two independent
reviews, using \cite[Proposition~4.3]{AS20Kleiman}.  The odd
weight--monodromy implication is recovered OW19, not a new
weight--monodromy theorem.  The ordinary-actual reconstruction,
finite trace, dominant detection, PH16 and GA17 remain frozen
dependencies.  R16 already proves arbitrary correspondence-kernel
covariance, abelian full G, supplied abelian-motive transfer,
nilpotence certificates, and the good-factor-times-abelian family.
The exact earlier candidate is R9,
\texttt{advance/obstruction\_good\_factor\_lowdim\_degenerations.tex},
especially (GS.2)--(GS.4) and its complete local admissibility theorem.
Its surface argument decomposed the incidence map into Tate blocks
and a Picard/Albanese block.  That source and the final R9 README
were extracted without modification and retained in the present
history-evidence directory.  The README labels the geometric claims
as internal research candidates.  R14's final report did not revise
those flags.  R16's continuation explicitly warns that historical
ideas are not activated merely by appearing in old files, while its
general numerical-model and dominant-detection bridges can of course
be applied to any independently verified earlier construction.
R17's release bridge is limited to its listed scope/status changes;
it is not a blanket certification of R9's product argument.
These facts do not justify declaring a new geometric domain merely
from the absence of its name in later registers.  Accordingly AS20
is recorded as an independent new proof and current-interface
verification, without a claim that the displayed family was outside
every previously established or logically implied domain.
It establishes this particular scope with the numerical-motive/curve-error
argument and the current frozen analytic bridge; it does not rely
on the earlier candidate proof.  It does not activate the separate
R9 multiple-bad-surface or Hilbert-scheme candidates.  No external
first-discovery claim is made for the present application.

\begin{thebibliography}{AS20}
\bibitem{AS20Jannsen}
U.~Jannsen, \emph{Motives, numerical equivalence, and semi-simplicity},
Invent. Math. \textbf{107} (1992), 447--452, Theorem~1.
\url{https://epub.uni-regensburg.de/26642/1/jannsen10.pdf}.
\bibitem{AS20Kleiman}
S.~L.~Kleiman, \emph{The standard conjectures}, in
\emph{Motives}, Proc. Sympos. Pure Math. \textbf{55}, Part~1
(1994), 3--20, Proposition~4.3.
\url{https://mishap.sdf.org/temp/Kleiman_Standard_Conjectures.pdf}.
\bibitem{AS20Saito}
T.~Saito, \emph{Weight spectral sequences and independence of $\ell$},
J. Inst. Math. Jussieu \textbf{2} (2003), 583--634,
Introduction, Conjecture~0.5 and the following known-case paragraph.
\url{https://www.ms.u-tokyo.ac.jp/~t-saito/pp/wmr2r.pdf}.
\end{thebibliography}
